\section{The row-local model on the primitive shell}
\label{sec:row-local-shell}

Throughout this section, $h\ne0$ is fixed.  Put
\begin{equation}
 \rho_h(a):=\one_{(a,h)=1}\frac{a}{\varphi(a)},
 \qquad a\geq1.
 \label{eq:row-local-density}
\end{equation}
This is the local von Mangoldt density on the row $ab+h$ after $b$ is
restricted by $(a,b)=1$.  If $(a,h)>1$, every value $ab+h$ has a fixed
prime divisor, apart from the exceptional possibility that it is a power
of that prime; the appropriate main density is therefore zero.  If
$(a,h)=1$, inclusion--exclusion over $(a,b)=1$ gives density
$a/\varphi(a)$ rather than the ambient density one.

Recall $u(a)=\Lambda(a)-1$.  Decompose the centered primitive shell as
\begin{align}
 \cK_h(X)
 &:={}
 \sum_{\substack{a>1,\ b\geq1\\(a,b)=1}}
 u(a)\bigl(\Lambda(ab+h)-1\bigr)W(ab/X),
 \label{eq:row-shell-K}\\
 \cM_h(X)
 &:={}
 \sum_{\substack{a>1,\ b\geq1\\(a,b)=1}}
 u(a)\bigl(\rho_h(a)-1\bigr)W(ab/X),
 \label{eq:row-shell-M}\\
 \cR_h(X)
 &:={}
 \sum_{\substack{a>1,\ b\geq1\\(a,b)=1}}
 u(a)\bigl(\Lambda(ab+h)-\rho_h(a)\bigr)W(ab/X).
 \label{eq:row-shell-R}
\end{align}
Thus, exactly at every finite $X$,
\begin{equation}
 \cK_h(X)=\cM_h(X)+\cR_h(X).
 \label{eq:row-shell-decomposition}
\end{equation}
Define
\begin{equation}
 \kappa_h:=\frac1{\zeta(2)}-\frac{\varphi(|h|)}{|h|},
 \qquad
 \delta_h:=\frac{\varphi(|h|)}{|h|}-\beta_h.
 \label{eq:kappa-delta-definitions}
\end{equation}
The two constants add to the coefficient in the fixed-shell theorem of
TPC-14:
\begin{equation}
 \kappa_h+\delta_h=\frac1{\zeta(2)}-\beta_h.
 \label{eq:kappa-delta-sum}
\end{equation}
Moreover,
\begin{equation}
 \beta_h=\frac{\varphi(|h|)}{|h|}
 \prod_{p\mid h}\left(1-\frac1{p^2-p+1}\right).
 \label{eq:beta-relative-to-phi}
\end{equation}
Hence $\delta_h>0$ whenever $|h|>1$, while
$\delta_{1}=\delta_{-1}=0$.

The main term of $\cM_h$ is not obtained by applying Euler summation to
each row and then summing the individual endpoint errors.  That procedure
is too wasteful when $a$ is as large as $X$.  We instead group by the
product and compute the relevant Dirichlet series.

For $r\geq1$, put
\begin{equation}
 m_h(r):=
 \sum_{\substack{a\parallel r\\a>1}}
 u(a)\bigl(\rho_h(a)-1\bigr),
 \label{eq:row-model-product-coefficient}
\end{equation}
where $a\parallel r$ means that $a\mid r$ and $(a,r/a)=1$.  Then
\begin{equation}
 \cM_h(X)=\sum_{r\geq1}m_h(r)W(r/X).
 \label{eq:row-model-product-grouping}
\end{equation}

\begin{lemma}[Dirichlet factorization of the row model]
\label{lem:row-model-dirichlet-factorization}
For $\Re s>1$,
\begin{equation}
 \sum_{r\geq1}\frac{m_h(r)}{r^s}
 =\zeta(s)^2\left(\frac1{\zeta(2s)}-C_h(s)\right)
 +\zeta(s)L_h(s),
 \label{eq:row-model-dirichlet-factorization}
\end{equation}
where
\begin{align}
 C_h(s)
 &:={}
 \prod_{p\mid h}(1-p^{-s})
 \prod_{p\nmid h}
 (1-p^{-s})
 \left(1+\frac{p^{1-s}}{p-1}\right),
 \label{eq:Ch-euler-product}\\
 L_h(s)
 &:={}
 \sum_{p\nmid h}\frac{\log p}{p-1}p^{-s}
 -\sum_{p\mid h}(\log p)p^{-s}.
 \label{eq:Lh-prime-series}
\end{align}
Both $C_h$ and $L_h$ are holomorphic in a neighborhood of $s=1$, and
\begin{equation}
 C_h(1)=\frac{\varphi(|h|)}{|h|}.
 \label{eq:Ch-at-one}
\end{equation}
Consequently, the double-pole coefficient at $s=1$ in
\eqref{eq:row-model-dirichlet-factorization} is $\kappa_h$.
\end{lemma}

\begin{proof}
For fixed $a$ one has
\[
 \sum_{\substack{b\geq1\\(a,b)=1}}b^{-s}
 =\zeta(s)\prod_{p\mid a}(1-p^{-s}).
\]
Since $\rho_h(1)=1$, the formally missing term $a=1$ contributes zero.
It follows from \eqref{eq:row-model-product-coefficient} that the left
side of \eqref{eq:row-model-dirichlet-factorization} equals
\begin{equation}
 \zeta(s)
 \sum_{a\geq1}
 \frac{(\Lambda(a)-1)(\rho_h(a)-1)}{a^s}
 \prod_{p\mid a}(1-p^{-s}).
 \label{eq:row-model-pre-factorization}
\end{equation}

We first treat the contribution of the constant part $-1$ in
$\Lambda-1$.  The two multiplicative series appearing there are
\begin{align*}
 \sum_{a\geq1}\frac1{a^s}\prod_{p\mid a}(1-p^{-s})
 &=\prod_p(1+p^{-s})=\frac{\zeta(s)}{\zeta(2s)},\\
 \sum_{a\geq1}\frac{\rho_h(a)}{a^s}
 \prod_{p\mid a}(1-p^{-s})
 &=\prod_{p\nmid h}
 \left(1+\frac{p^{1-s}}{p-1}\right)
 =\zeta(s)C_h(s).
\end{align*}
For $p\nmid h$, the corresponding local factor in $C_h$ is
\[
 (1-p^{-s})\left(1+\frac{p^{1-s}}{p-1}\right)
 =1+\frac{p^{-s}}{p-1}-\frac{p^{1-2s}}{p-1}.
\]
The product therefore converges normally on compact subsets of a
half-plane containing $s=1$.  At $s=1$ this local factor is exactly one.
Only the finitely many primes dividing $h$ remain, which proves
\eqref{eq:Ch-at-one}.

For the $\Lambda$ part, only prime powers occur.  If $p\nmid h$, their
total contribution to the inner series in
\eqref{eq:row-model-pre-factorization} is
\[
 \sum_{j\geq1}
 (\log p)\left(\frac{p}{p-1}-1\right)
 (1-p^{-s})p^{-js}
 =\frac{\log p}{p-1}p^{-s}.
\]
If $p\mid h$, it is instead $-(\log p)p^{-s}$.  This gives
\eqref{eq:Lh-prime-series}.  The first prime sum is normally convergent
for $\Re s>0$, locally around $s=1$, and the second is finite.  Multiplying
the resulting inner identity by the outer $\zeta(s)$ in
\eqref{eq:row-model-pre-factorization} proves
\eqref{eq:row-model-dirichlet-factorization}.  Its only double pole at
$s=1$ comes from the first term, with coefficient
$1/\zeta(2)-C_h(1)=\kappa_h$.
\end{proof}

\begin{theorem}[Global row-model law]
\label{thm:global-row-model-law}
For every fixed nonzero $h$,
\begin{equation}
 \cM_h(X)
 =\kappa_hXI_0(W)\log X+O_{h,W}(X).
 \label{eq:global-row-model-law}
\end{equation}
\end{theorem}

\begin{proof}
The Euler products and prime series in
\cref{lem:row-model-dirichlet-factorization} converge normally for
$\Re s>1/2$ and $\Re s>0$, respectively.  Since $1/\zeta(2s)$ is
holomorphic for $\Re s>1/2$, the factorization continues meromorphically
to that half-plane, with its only pole on $\Re s\geq1$ at $s=1$.  It is a
double pole with leading coefficient $\kappa_h$; all remaining terms have
at most a simple pole there.

Let $\widehat W(s)=\int_0^\infty W(x)x^{s-1}\,dx$.  Mellin inversion on
$\Re s>1$, followed by a shift to $\Re s=1/2+\varepsilon$, is legitimate:
$\widehat W$ decays faster than every power in vertical strips, while the
other factors have polynomial vertical growth.  The residue at $s=1$
therefore gives
\[
 \sum_rm_h(r)W(r/X)
 =\kappa_hXI_0(W)\log X+c_{h,W}X+o_{h,W}(X)
\]
 for a finite secondary constant $c_{h,W}$, with an error
$O_{h,W,\varepsilon}(X^{1/2+\varepsilon})$.  Absorbing the secondary term
and using \eqref{eq:row-model-product-grouping} proves the theorem.
\end{proof}

TPC-14 \citep{WangTPC14} evaluates the original shell, for the same fixed
$h$, as
\begin{equation}
 \cK_h(X)
 =\left(\frac1{\zeta(2)}-\beta_h\right)
 XI_0(W)\log X+O_{h,W}(X).
 \label{eq:imported-fixed-shell-law}
\end{equation}
Subtracting the row model isolates what is not explained by local row
density.

\begin{corollary}[Global row-residual law]
\label{cor:global-row-residual-law}
For every fixed nonzero $h$,
\begin{equation}
 \cR_h(X)
 =\delta_hXI_0(W)\log X+O_{h,W}(X).
 \label{eq:global-row-residual-law}
\end{equation}
In particular,
\begin{equation}
 \kappa_2=\frac6{\pi^2}-\frac12,
 \qquad
 \delta_2=\frac16.
 \label{eq:h-two-kappa-delta}
\end{equation}
\end{corollary}

\begin{proof}
Use the exact identity \eqref{eq:row-shell-decomposition}, then subtract
\eqref{eq:global-row-model-law} from
\eqref{eq:imported-fixed-shell-law}.  The coefficient is
\[
 \left(\frac1{\zeta(2)}-\beta_h\right)-\kappa_h
 =\frac{\varphi(|h|)}{|h|}-\beta_h=\delta_h.
\]
For $h=2$, one has $\varphi(2)/2=1/2$ and $\beta_2=1/3$.
\end{proof}

\begin{remark}[What the residual law says]
\label{rem:global-row-residual-interpretation}
The quantity $\cR_h$ is still summed over every primitive orientation
$b$.  Equation \eqref{eq:global-row-residual-law} says that local
row-density calibration does not explain the full fixed-shell main term.
It neither assigns the remaining term to any individual orientation nor
gives information about the endpoint $b=1$.
\end{remark}
